\maketitle
\head{本版编排与收录说明}
\noindent\textbf{本书是当前已译内容的中文工作合订稿，尚非全部上传材料的完结全译本。}
第十九版以已发布的第十八版实际源码为基础，继续对读第五卷预备笔记第一篇的8.14至第11节末：法文物理46页下部至63页，英文32至43页第12节以前。旧译按段替换，作者证明、脚注、删划与编辑注继续保存。本轮修复固定分支和维数条件等实质问题，也恢复英译漏失的法稿条件；译者补证均另标身份。此前已收入的Dot原批号69--108及所有其他已译正文保留。

\section*{正文次序}
依原著内容排列：导言；归并后的第0章预备知识；第I章概形的语言；第II章若干类态射的初等整体研究；第III章凝聚层的上同调研究；第IV章待译范围说明；第V卷预备笔记；各分册书末资料和来源记录。第III章两部分的§§1--5与§§6--7直接衔接。第0章原分载于第一卷、第三卷第一部分及第四卷第一部分；保留其出处和原著交叉引用。

第V卷沿用既有按内容排列的次序：先列函数的奇异零点与超奇异零点、光滑型和初等奇点，再接光滑二次型附录；随后为超平面截面与锥投影第一篇，之后是第三篇射影丛、相对除子与线性系，最后收入致迪厄多内的书信、手写提纲补遗及三块删划草稿的补译。PS短篇与合并英文稿的同文只排印一次。法文扫描的物理页序与手稿阅读次序不完全一致，具体对应见译注及工程来源表。

\clearpage
\section*{实际收录范围}
\begin{longtable}{@{}p{.25\linewidth}p{.67\linewidth}@{}}
\toprule 材料&本版源码中的连续范围及状态\\\midrule\endhead
第一卷（1960）&导言；第0章§§1--7；第I章§§1--10，止于（10.15.7）及证明；相关书末资料。正文沿用既有源码。\\
第二卷（1961）&第II章§§1--8，止于（8.14.14）及证明；参考文献、索引、目录及勘误清单一。正文沿用既有源码。\\
第三卷第一部分（1961）&第0章§§8--13；第III章§§1--5，止于（5.5.2）；相关书末资料。\\
第三卷第二部分（1963）&第III章§§6--7，止于（7.9.14）；记号、术语索引、原目录及勘误清单二。\\
第四卷第一部分（1964）&第0章补篇§§14--23，止于（23.2.10，ii）及原“待续”。正式第IV章主体仍未连续译入。\\
第V卷预备笔记&既有第一篇§§1--15和短篇及其后续；既有英文物理78--112页的光滑二次型附录、第三篇及书信，以及法文物理150页提纲和153页异扫说明均保留。第一篇§§1--7及8.1--8.13沿用第十八版回校稿；本版继续回校8.14至第11节，§§12--15沿用旧稿。原稿缺文、删划与未决事项仍明示。\\
\bottomrule
\end{longtable}

\section*{本版整合与订正的性质}
本轮续校针对原稿和现存译文，不把已经收入的Dot译稿重新计为新译。主要改动包括：8.16与8.18原条件的反例及明确加强范围；一般线性系的固定分支和投影维数；9.3的附加条件；爆破条目的正则性范围；第11节法稿中被英译遗漏的条件。每项错误、原说法与订正依据见正文脚注及累计核查报告的新篇。旧第十八版工程与报告全部另行保存。

已确认的原文错误采用有依据的订正，原错及理由写入脚注；缺文、难辨字和未证的一般化保持标记。原作者文字、原稿删划、英文编者说明、原批译者补证及本轮整合说明各按其身份阅读。第8.7(c)的缺续等问题没有因更新版次而消失；对某一附加假设下的补证也不能推广为整条原断言已经成立。

本版没有重新逐句审定第I、II、III章和第0章，也没有把全部历史报告的意见重新证明一遍。本轮实际的来源对读、改动、编译和版面检查范围，见累计报告开头及工程检查记录。原报告的“本轮”均指其制作当时。

\section*{真实接续位置}
第五卷英文底本的后续材料已经接到其物理末页112页的书信；法文扫描另有物理150页提纲和153页异扫的处理记录。到达英文末页不表示法文171页全部完成逐页校勘，也不把不存在的英文第113页当作续译点。

第一篇本版法英回校止于第11节末（法文物理63页；英文物理43页第12节以前）。下一回校位置为法文物理64页的第12节，对读英文物理43页中段；既有对应源码为\texttt{continuations/v\_dot\_middle.tex}。这些后文已经有译文，下一步是据底本严谨回校。

正式第IV章仍另从\texttt{EGAIV.1.pdf}物理221页内容提要、222--223页章首说明及223页（1.1.1）接续。第0章§§14--23与预备手稿的拟定章号均不能替代正式第IV章主体。

\section*{版式与使用}
使用真正的\texttt{\string\documentclass[UTF8,a4paper]\{ctexbook\}}。字体、字号、正文版心、标题和双面分页由类文件默认决定，不另设字体或页边距。原章号、条目号和公式号保留，交换图以TeX排版。总目录指向本合订本；原分册目录所列页码仍指原版。为四位页码保留目录宽度，正文版心不变。默认双面分页所需空白页保留。

提供合订PDF、可独立编译的单文件TeX、累计核查报告和工程。单文件无需其他正文输入，仍需XeLaTeX、相关宏包和默认字体。源码清单、构建日志、逐页比较和抽样渲染各有具体用途，均不代替全书数学审定。
\tableofcontents
